\begin{abstract}
本稿では，画像中の物体らしい領域を検出するタスクと，その領域のカテゴリをアノテーションするタスクを分離し，WordNetの下位語をたどることで追加学習なしに粒度を変更できる物体アノテーション手法を検討する．
提案手法はクラス非依存の検出にDETRを用い，各領域のクロップのカテゴリをSigLIPのゼロショット分類で推定する．
ハイブリッド手法はYOLOの予測を深さ0とし，YOLOのボックスのクロップに対してSigLIPとWordNetで粒度を変更する．
さらに，YOLOが検出しないCOCO外の物体を，物体らしさに基づく検出器の候補で補う構成を提案する．
COCO 2017 ValidationとLVISのアノテーションを用い，検証用とテスト用に分けた画像で評価した結果，以下を確認した．
第一に，COCOの80クラスの範囲では，ハイブリッド手法がYOLO11xと同じTop-1正解率0.969を保ちつつ，粒度を変更できたインスタンスを提案手法の1.79倍に増やした．
第二に，LVISの細かい粒度の正解に対し，シグモイドの累積確信度で細分化を止めるハイブリッド手法の階層的F値は0.659となり，YOLO11x単体の0.604を上回った．
第三に，COCOで学習したYOLO11xとDETRは閾値を下げてもCOCO外の物体の再現率が2割未満で飽和する一方，OWLv2の物体らしさのスコアは1画像あたり100件で74.3\%に達した．
第四に，YOLO11xのボックスと重ならないOWLv2の候補を加えることで，COCOの物体の精度を保ったまま，COCO外の物体の検出を4.6\%から31.6\%に増やし，粒度の変更の対象にできた．
ただし，検出したCOCO外の物体のうち深さ0を正しく命名できたのは8.1\%であり，命名の語彙と不適切な領域の棄却に課題が残る．
\add{そこで，物体らしさとシグモイドの確信度による候補の棄却と，頻出するカテゴリの語から順に下位語を選ぶ命名，およびSigLIP SO400Mを組み合わせた結果，1画像あたり5件の候補を加える条件で，COCO外の物体の命名を1.7\%から4.1\%に，加えた候補の命名の適合率を3.6\%から8.1\%に向上させた．}
\end{abstract}

\noindent
{\bf キーワード：}物体アノテーション，粒度可変アノテーション，WordNet，SigLIP，YOLO，OWLv2，オープンボキャブラリ
